\documentclass{article}
\usepackage{mathtools} 
\usepackage{fontspec}
\usepackage[UTF8]{ctex}
\usepackage{amsthm}
\usepackage{mdframed}
\usepackage{xcolor}
\usepackage{amssymb}
\usepackage{amsmath}


% 定义新的带灰色背景的说明环境 zremark
\newmdtheoremenv[
  backgroundcolor=gray!10,
  % 边框与背景一致，边框线会消失
  linecolor=gray!10
]{zremark}{说明}

% 通用矩阵命令: \flexmatrix{矩阵名}{元素符号}{行数}{列数}
\newcommand{\flexmatrix}[4]{
  \[
  #1 = \begin{pmatrix}
    #2_{11}     & #2_{12}     & \cdots & #2_{1#4}   \\
    #2_{21}     & #2_{22}     & \cdots & #2_{2#4}   \\
    \vdots      & \vdots      & \ddots & \vdots     \\
    #2_{#31}    & #2_{#32}    & \cdots & #2_{#3#4}
  \end{pmatrix}
  \]
}

% 简化版命令（默认矩阵名为A，元素符号为a）: \quickmatrix{行数}{列数}
\newcommand{\quickmatrix}[2]{\flexmatrix{A}{a}{#1}{#2}}

\begin{document}
\title{注释 17.2}
\author{张志聪}
\maketitle

\begin{zremark}
  定理17.2（可微的充分条件）的证明中为什么在$f(x_0 + \Delta x, y_0 + \Delta y) - f(x_0, y_0 + \Delta y)$
  上可以使用拉格朗日定理，定理的条件是如何成立的，尤其是$y = y_0 + \Delta y, x \in (x_0, x_0 + \Delta x)$
  偏导数$f_x(x, y)$为什么存在。
\end{zremark}

\textbf{证明：}
这里，只需要对$f_x$在$(x_0, y_0)$处连续有正确的理解，即可。

要把$f_x(x, y)$看做二元函数，$f_x$在$(x_0, y_0)$处连续，也就是说：
$(x, y)$以任意路径趋近$(x_0, y_0)$，$f_x(x, y)$的值都趋近于同一个数。
等价于
\begin{align*}
  \lim\limits_{(x, y) \to (x_0, y_0)} f_x(x, y) = f_x(x_0, y_0)
\end{align*}
所以存在某个邻域$U(P_0)$，$f_x(x, y)$在$U(P_0)$有定义。
因为$(\Delta x, \Delta y) \to (0, 0)$，故讨论的点都在邻域$U(P_0)$内，
所以$y = y_0 + \Delta y, x \in (x_0, x_0 + \Delta x)$
偏导数$f_x(x, y)$都存在。

一种错误的理解：
\begin{align*}
  \lim\limits_{x \to x_0} f_x(x, y_0) = f_x(x_0, y_0)
\end{align*}
即把$y$看做固定值$y_0$，这是把连续性和偏导数在一点处的计算方式混淆在一起了。




\end{document}